\section{Permutation representations and transitive actions}

Let
\[
    \mathcal{C}=\operatorname{Set}.
\]

An object of \(\operatorname{Rep}(G,\operatorname{Set})\) is a pair
\[
    (X,\rho),
\]
where \(X\) is a set and
\[
    \rho:G\longrightarrow \operatorname{Sym}(X)
\]
is a group homomorphism. Equivalently, \(G\) acts on \(X\) by
\[
    g\cdot x:=\rho(g)(x),
    \qquad g\in G,\ x\in X.
\]

A morphism
\[
    f:(X,\rho)\longrightarrow (Y,\sigma)
\]
is a map of sets
\[
    f:X\longrightarrow Y
\]
such that, for every \(g\in G\), the following diagram commutes:
\[
\begin{tikzcd}
    X \arrow[r,"f"] \arrow[d,"\rho(g)"']
        & Y \arrow[d,"\sigma(g)"] \\
    X \arrow[r,"f"']
        & Y
\end{tikzcd}
\]
Equivalently,
\[
    f(\rho(g)(x))=\sigma(g)(f(x)),
    \qquad g\in G,\ x\in X.
\]
In action notation, this condition is written as
\[
    f(g\cdot x)=g\cdot f(x).
\]

\begin{lemma}
    Let
    \[
        f:(X,\rho)\longrightarrow (Y,\sigma)
    \]
    be a morphism in \(\operatorname{Rep}(G,\operatorname{Set})\). If the
    underlying map \(f:X\to Y\) is bijective, then \(f\) is an isomorphism in
    \(\operatorname{Rep}(G,\operatorname{Set})\).
\end{lemma}

\begin{proof}
    Since \(f:X\to Y\) is bijective, it has an inverse map
    \[
        f^{-1}:Y\longrightarrow X.
    \]
    We need to show that \(f^{-1}\) is \(G\)-equivariant.

    Let \(y\in Y\). Write \(y=f(x)\) for a unique \(x\in X\). For every \(g\in G\),
    since \(f\) is equivariant, we have
    \[
        f(g\cdot x)=g\cdot f(x)=g\cdot y.
    \]
    Applying \(f^{-1}\), we get
    \[
        f^{-1}(g\cdot y)=g\cdot x=g\cdot f^{-1}(y).
    \]
    Hence \(f^{-1}\) is \(G\)-equivariant. Therefore \(f\) is an isomorphism in
    \(\operatorname{Rep}(G,\operatorname{Set})\).
\end{proof}

For sets \(X\) and \(Y\), recall the product
\[
    X\times Y
\]
and the disjoint union
\[
    X\sqcup Y.
\]
If \(Z\) is another set, then there is a natural bijection
\[
    Z\times (X\sqcup Y)
    \cong
    (Z\times X)\sqcup (Z\times Y).
\]

Let \((X,\rho)\) and \((Y,\sigma)\) be two \(G\)-sets.

\begin{definition}
    The disjoint union of \((X,\rho)\) and \((Y,\sigma)\) is the \(G\)-set
    \[
        (X,\rho)\sqcup (Y,\sigma)
        :=
        (X\sqcup Y,\rho\sqcup\sigma),
    \]
    where the action is defined by
    \[
        (\rho\sqcup\sigma)(g)(z)
        =
        \begin{cases}
            \rho(g)(z), & z\in X,\\
            \sigma(g)(z), & z\in Y.
        \end{cases}
    \]
\end{definition}

\begin{definition}
    The product of \((X,\rho)\) and \((Y,\sigma)\) is the \(G\)-set
    \[
        (X,\rho)\times (Y,\sigma)
        :=
        (X\times Y,\rho\times\sigma),
    \]
    where
    \[
        (\rho\times\sigma)(g)(x,y)
        =
        \bigl(\rho(g)(x),\sigma(g)(y)\bigr).
    \]
    In action notation,
    \[
        g\cdot (x,y)=(g\cdot x,g\cdot y).
    \]
\end{definition}

\begin{definition}
    Let \((X,\rho)\) be a set-representation of \(G\). A subset
    \[
        X_1\subseteq X
    \]
    is called \(G\)-stable if
    \[
        \rho(g)(X_1)\subseteq X_1
    \]
    for every \(g\in G\).

    If \(X_1\subseteq X\) is \(G\)-stable, then the restricted action
    \[
        \rho_1(g):=\rho(g)|_{X_1}
    \]
    defines a representation
    \[
        (X_1,\rho_1),
    \]
    called a subrepresentation of \((X,\rho)\).
\end{definition}

\begin{remark}
    Since every \(\rho(g)\) is a bijection and \(\rho(g^{-1})=\rho(g)^{-1}\), the
    condition
    \[
        \rho(g)(X_1)\subseteq X_1
    \]
    for every \(g\in G\) implies
    \[
        \rho(g)(X_1)=X_1
    \]
    for every \(g\in G\).
\end{remark}

\begin{example}
    Let \((X,\rho)\) be a \(G\)-set and let \(x\in X\). The orbit of \(x\) is
    \[
        X_1:=Gx:=\{\,g\cdot x\mid g\in G\,\}.
    \]
    Then \(X_1\) is \(G\)-stable. Indeed, for \(h\in G\),
    \[
        h\cdot (g\cdot x)=(hg)\cdot x\in Gx.
    \]
    Therefore
    \[
        (X_1,\rho|_{X_1})
    \]
    is a subrepresentation of \((X,\rho)\).
\end{example}

\begin{proposition}
    Let \((X,\rho)\) be a \(G\)-set and let
    \[
        X_1=Gx
    \]
    be the orbit of some \(x\in X\). Put
    \[
        X_2:=X\setminus X_1.
    \]
    Then \(X_2\) is also \(G\)-stable. Hence
    \[
        (X,\rho)
        =
        (X_1,\rho_1)\sqcup (X_2,\rho_2),
    \]
    where
    \[
        \rho_1=\rho|_{X_1},
        \qquad
        \rho_2=\rho|_{X_2}.
    \]
\end{proposition}

\begin{proof}
    We already know that \(X_1\) is \(G\)-stable. We prove that \(X_2\) is
    \(G\)-stable.

    Let \(y\in X_2\) and \(g\in G\). Suppose, for contradiction, that
    \[
        g\cdot y\in X_1.
    \]
    Since \(X_1\) is \(G\)-stable and \(g^{-1}\in G\), applying \(g^{-1}\) gives
    \[
        y=g^{-1}\cdot (g\cdot y)\in X_1,
    \]
    contradicting \(y\in X_2=X\setminus X_1\). Hence
    \[
        g\cdot y\in X_2.
    \]
    Therefore \(X_2\) is \(G\)-stable.

    Since
    \[
        X=X_1\sqcup X_2
    \]
    as sets and both \(X_1\) and \(X_2\) are \(G\)-stable, the action \(\rho\)
    restricts to actions \(\rho_1\) and \(\rho_2\), and we obtain
    \[
        (X,\rho)
        =
        (X_1,\rho_1)\sqcup (X_2,\rho_2).
    \]
\end{proof}

Let \((X,\rho)\) be a set-representation of \(G\), that is,
\[
    \rho:G\longrightarrow \operatorname{Sym}(X).
\]
We write
\[
    g\cdot x:=\rho(g)(x),
    \qquad g\in G,\ x\in X.
\]

\begin{proposition}
    Let \((X,\rho)\) be a set-representation of \(G\), and assume that
    \[
        X\neq \varnothing.
    \]
    The following conditions are equivalent:
    \begin{enumerate}
        \item \((X,\rho)\) is simple, or irreducible. That is, if
        \((X_1,\rho_1)\) is a subrepresentation of \((X,\rho)\), then
        \[
            X_1=X
            \qquad\text{or}\qquad
            X_1=\varnothing.
        \]

        \item \((X,\rho)\) is indecomposable. That is, if
        \[
            (X,\rho)
            =
            (X_1,\rho_1)\sqcup (X_2,\rho_2),
        \]
        then
        \[
            X_1=\varnothing
            \qquad\text{or}\qquad
            X_2=\varnothing.
        \]

        \item \((X,\rho)\) is transitive. That is, for every \(x,y\in X\), there
        exists \(g\in G\) such that
        \[
            y=g\cdot x.
        \]
    \end{enumerate}
\end{proposition}

\begin{proof}
    We prove the equivalence of the three conditions.

    \textup{(1) \(\Rightarrow\) (2).}
    Suppose that \((X,\rho)\) is simple and that
    \[
        (X,\rho)
        =
        (X_1,\rho_1)\sqcup (X_2,\rho_2).
    \]
    Then \(X_1\subseteq X\) is a \(G\)-stable subset, hence
    \((X_1,\rho_1)\) is a subrepresentation of \((X,\rho)\). By simplicity,
    \[
        X_1=X
        \qquad\text{or}\qquad
        X_1=\varnothing.
    \]
    If \(X_1=X\), then \(X_2=\varnothing\). Therefore \((X,\rho)\) is
    indecomposable.

    \textup{(2) \(\Rightarrow\) (1).}
    Suppose that \((X,\rho)\) is indecomposable. Let \(X_1\subseteq X\) be a
    \(G\)-stable subset. Then its complement
    \[
        X_2:=X\setminus X_1
    \]
    is also \(G\)-stable. Hence
    \[
        (X,\rho)
        =
        (X_1,\rho|_{X_1})\sqcup (X_2,\rho|_{X_2}).
    \]
    By indecomposability,
    \[
        X_1=\varnothing
        \qquad\text{or}\qquad
        X_2=\varnothing.
    \]
    In the second case \(X_1=X\). Thus \((X,\rho)\) is simple.

    \textup{(1) \(\Leftrightarrow\) (3).}
    For \(x\in X\), the orbit
    \[
        Gx:=\{\,g\cdot x\mid g\in G\,\}
    \]
    is a \(G\)-stable subset of \(X\). If \((X,\rho)\) is simple, then since
    \(Gx\neq \varnothing\), we must have
    \[
        Gx=X.
    \]
    Thus every element of \(X\) lies in the orbit of \(x\), so the action is
    transitive.

    Conversely, if the action is transitive and \(X_1\subseteq X\) is a nonempty
    \(G\)-stable subset, choose \(x\in X_1\). For any \(y\in X\), transitivity gives
    some \(g\in G\) such that
    \[
        y=g\cdot x.
    \]
    Since \(X_1\) is \(G\)-stable, \(g\cdot x\in X_1\). Hence \(y\in X_1\), so
    \(X_1=X\). Therefore there are no nontrivial subrepresentations, and
    \((X,\rho)\) is simple.
\end{proof}

\begin{example}[The transitive \(G\)-set \(G/H\)]
    Let \(H\leq G\). Consider the set of left cosets
    \[
        G/H:=\{\,xH\mid x\in G\,\}.
    \]
    Define an action of \(G\) on \(G/H\) by
    \[
        \rho_H(g)(xH)=gxH.
    \]
    This action is transitive. Indeed, for any two cosets \(xH,yH\in G/H\), we have
    \[
        yH=(yx^{-1})xH.
    \]
    Hence
    \[
        (G/H,\rho_H)
    \]
    is a simple set-representation of \(G\).

    Moreover, the stabilizer of the coset \(yH\) is
    \[
        G_{yH}=yHy^{-1}.
    \]
    Indeed,
    \[
        g\in G_{yH}
        \iff
        gyH=yH
        \iff
        y^{-1}gy\in H
        \iff
        g\in yHy^{-1}.
    \]
\end{example}

Let \((X,\rho)\) be a \(G\)-set. Define a relation \(\sim_G\) on \(X\) by
\[
    x\sim_G y
    \iff
    \exists g\in G \text{ such that } y=g\cdot x.
\]
This is an equivalence relation. Its equivalence classes are precisely the orbits of
the \(G\)-action.

We denote by
\[
    G\backslash X
\]
the set of equivalence classes, or equivalently, the set of \(G\)-orbits in \(X\).

\begin{proposition}
    Let \((X,\rho)\) be a \(G\)-set.
    \begin{enumerate}
        \item Any equivalence class \(\Omega\in G\backslash X\) defines a simple
        subrepresentation
        \[
            (\Omega,\rho|_{\Omega}).
        \]

        \item There is a decomposition
        \[
            (X,\rho)
            =
            \bigsqcup_{\Omega\in G\backslash X}
            (\Omega,\rho|_{\Omega}).
        \]
    \end{enumerate}
\end{proposition}

\begin{proof}
    Each equivalence class \(\Omega\in G\backslash X\) is an orbit. Hence it is
    \(G\)-stable, and the restriction
    \[
        \rho|_{\Omega}:G\longrightarrow \operatorname{Sym}(\Omega)
    \]
    defines a subrepresentation \((\Omega,\rho|_{\Omega})\).

    The action of \(G\) on \(\Omega\) is transitive by definition of an orbit.
    Therefore, by the previous proposition, \((\Omega,\rho|_{\Omega})\) is simple.

    Finally, the set \(X\) is the disjoint union of its equivalence classes:
    \[
        X=\bigsqcup_{\Omega\in G\backslash X}\Omega.
    \]
    Since each \(\Omega\) is \(G\)-stable, this gives a decomposition of
    \(G\)-sets:
    \[
        (X,\rho)
        =
        \bigsqcup_{\Omega\in G\backslash X}
        (\Omega,\rho|_{\Omega}).
    \]
\end{proof}

\begin{theorem}
    Let \((X,\rho)\) be a transitive set-representation of \(G\). For \(x\in X\),
    let
    \[
        G_x:=\{\,g\in G\mid g\cdot x=x\,\}
    \]
    be the stabilizer of \(x\). Then the map
    \[
        \Phi:G/G_x\longrightarrow X,
        \qquad
        gG_x\longmapsto g\cdot x
    \]
    defines an isomorphism of \(G\)-sets
    \[
        (G/G_x,\rho_{G_x})
        \cong
        (X,\rho).
    \]
\end{theorem}

\begin{proof}
    First we show that \(\Phi\) is well-defined. Suppose
    \[
        gG_x=g'G_x.
    \]
    Then \(g'=gh\) for some \(h\in G_x\). Hence
    \[
        g'\cdot x=(gh)\cdot x=g\cdot (h\cdot x)=g\cdot x,
    \]
    since \(h\in G_x\). Thus \(\Phi(gG_x)=\Phi(g'G_x)\).

    The map \(\Phi\) is surjective because the action of \(G\) on \(X\) is
    transitive.

    It is injective as well. If
    \[
        \Phi(gG_x)=\Phi(g'G_x),
    \]
    then
    \[
        g\cdot x=g'\cdot x.
    \]
    Hence
    \[
        (g'^{-1}g)\cdot x=x,
    \]
    so \(g'^{-1}g\in G_x\). Therefore
    \[
        gG_x=g'G_x.
    \]

    Finally, \(\Phi\) is \(G\)-equivariant. For \(a\in G\),
    \[
        \Phi\bigl(\rho_{G_x}(a)(gG_x)\bigr)
        =
        \Phi(agG_x)
        =
        ag\cdot x,
    \]
    while
    \[
        \rho(a)\bigl(\Phi(gG_x)\bigr)
        =
        a\cdot (g\cdot x)
        =
        ag\cdot x.
    \]
    Hence
    \[
        \Phi\circ \rho_{G_x}(a)
        =
        \rho(a)\circ \Phi.
    \]
    Thus \(\Phi\) is an isomorphism of \(G\)-sets.
\end{proof}

\begin{theorem}
    Let \(H,H'\leq G\). Then the \(G\)-sets
    \[
        (G/H,\rho_H)
        \qquad\text{and}\qquad
        (G/H',\rho_{H'})
    \]
    are isomorphic if and only if \(H\) and \(H'\) are conjugate in \(G\).
\end{theorem}

\begin{proof}
    Suppose first that
    \[
        F:G/H\longrightarrow G/H'
    \]
    is an isomorphism of \(G\)-sets. Write
    \[
        F(H)=aH'
    \]
    for some \(a\in G\). Since \(F\) is \(G\)-equivariant, the stabilizer of \(H\)
    must equal the stabilizer of \(F(H)\). Now
    \[
        G_H=H,
    \]
    and
    \[
        G_{aH'}=aH'a^{-1}.
    \]
    Hence
    \[
        H=aH'a^{-1}.
    \]
    Thus \(H\) and \(H'\) are conjugate in \(G\).

    Conversely, suppose that \(H\) and \(H'\) are conjugate. Write
    \[
        H=aH'a^{-1}
    \]
    for some \(a\in G\). Define
    \[
        F:G/H\longrightarrow G/H'
    \]
    by
    \[
        F(gH)=gaH'.
    \]
    This is well-defined. Indeed, if \(gH=g_1H\), then \(g_1=gh\) for some
    \(h\in H\). Since \(H=aH'a^{-1}\), we have
    \[
        a^{-1}ha\in H'.
    \]
    Therefore
    \[
        g_1aH'=ghaH'=ga(a^{-1}ha)H'=gaH'.
    \]
    Thus \(F(gH)=F(g_1H)\).

    The map \(F\) is bijective, with inverse
    \[
        F^{-1}(gH')=ga^{-1}H.
    \]
    It is also \(G\)-equivariant, since for \(t\in G\),
    \[
        F(tgH)=tgaH'=tF(gH).
    \]
    Hence \(F\) is an isomorphism of \(G\)-sets:
    \[
        (G/H,\rho_H)\cong (G/H',\rho_{H'}).
    \]
\end{proof}

\begin{definition}[Degree]
    Let \((X,\rho)\) be a set-representation of \(G\), that is,
    \[
        \rho:G\longrightarrow \operatorname{Sym}(X).
    \]
    The degree of \((X,\rho)\) is defined to be the cardinality of \(X\), denoted by
    \[
        \deg(X,\rho):=|X|.
    \]
\end{definition}

\begin{corollary}
    Let \(G\) be a finite group. Then the degree of any irreducible
    set-representation of \(G\) divides \(|G|\).
\end{corollary}

\begin{proof}
    Let \((X,\rho)\) be an irreducible set-representation of \(G\). Irreducible set-representations are transitive \(G\)-sets, the action
    of \(G\) on \(X\) is transitive.

    Choose \(x\in X\), and let
    \[
        G_x:=\{\,g\in G\mid g\cdot x=x\,\}
    \]
    be the stabilizer of \(x\). By the classification of transitive \(G\)-sets, we
    have an isomorphism of \(G\)-sets
    \[
        X\cong G/G_x.
    \]
    Hence
    \[
        |X|=|G/G_x|=[G:G_x].
    \]
    Since \(G\) is finite, Lagrange's theorem gives
    \[
        |G|=[G:G_x]\cdot |G_x|.
    \]
    Therefore
    \[
        |X|=[G:G_x]
    \]
    divides \(|G|\). Thus the degree of \((X,\rho)\) divides \(|G|\).
\end{proof}

